
\begin{code}[hide]%
\>[0]\AgdaSymbol{\{-\#}\AgdaSpace{}%
\AgdaKeyword{OPTIONS}\AgdaSpace{}%
\AgdaPragma{--prop}\AgdaSpace{}%
\AgdaSymbol{\#-\}}\<%
\\
%
\\[\AgdaEmptyExtraSkip]%
\>[0]\AgdaKeyword{open}\AgdaSpace{}%
\AgdaKeyword{import}\AgdaSpace{}%
\AgdaModule{Agda.Primitive}\<%
\\
\>[0]\AgdaKeyword{open}\AgdaSpace{}%
\AgdaKeyword{import}\AgdaSpace{}%
\AgdaModule{Lib}\<%
\\
%
\\[\AgdaEmptyExtraSkip]%
\>[0]\AgdaKeyword{module}\AgdaSpace{}%
\AgdaModule{model}\<%
\\
\>[0][@{}l@{\AgdaIndent{0}}]%
\>[2]\AgdaSymbol{(}\AgdaBound{funar}\AgdaSpace{}%
\AgdaSymbol{:}\AgdaSpace{}%
\AgdaDatatype{ℕ}\AgdaSpace{}%
\AgdaSymbol{→}\AgdaSpace{}%
\AgdaPrimitive{Set}\AgdaSymbol{)}\<%
\\
%
\>[2]\AgdaSymbol{(}\AgdaBound{relar}\AgdaSpace{}%
\AgdaSymbol{:}\AgdaSpace{}%
\AgdaDatatype{ℕ}\AgdaSpace{}%
\AgdaSymbol{→}\AgdaSpace{}%
\AgdaPrimitive{Set}\AgdaSymbol{)}\<%
\\
%
\>[2]\AgdaKeyword{where}\<%
\\
%
\\[\AgdaEmptyExtraSkip]%
\>[0]\AgdaKeyword{record}\AgdaSpace{}%
\AgdaRecord{Model}\AgdaSpace{}%
\AgdaSymbol{\{}\AgdaBound{i}\AgdaSymbol{\}\{}\AgdaBound{j}\AgdaSymbol{\}\{}\AgdaBound{k}\AgdaSymbol{\}\{}\AgdaBound{l}\AgdaSymbol{\}}\AgdaSpace{}%
\AgdaSymbol{:}\AgdaSpace{}%
\AgdaPrimitive{Set}\AgdaSpace{}%
\AgdaSymbol{(}\AgdaPrimitive{lsuc}\AgdaSpace{}%
\AgdaSymbol{(}\AgdaBound{i}\AgdaSpace{}%
\AgdaOperator{\AgdaPrimitive{⊔}}\AgdaSpace{}%
\AgdaBound{j}\AgdaSpace{}%
\AgdaOperator{\AgdaPrimitive{⊔}}\AgdaSpace{}%
\AgdaBound{k}\AgdaSpace{}%
\AgdaOperator{\AgdaPrimitive{⊔}}\AgdaSpace{}%
\AgdaBound{l}\AgdaSymbol{))}\AgdaSpace{}%
\AgdaKeyword{where}\<%
\\
\>[0][@{}l@{\AgdaIndent{0}}]%
\>[3]\AgdaKeyword{field}\<%
\end{code}

\section{Kategória}

A modell fogalom első elemeként megadunk egy kontextusok és behelyettesítések által alkotott kategóriát. A kategória objektumait a kontextusok reprezentálják és a \AgdaField{Con} szort határozza meg. A morfizmusai pedig a behelyettesítések lesznek, amelyeket a \AgdaField{Sub} szort ad meg.
\begin{code}%
\>[3][@{}l@{\AgdaIndent{1}}]%
\>[6]\AgdaField{Con}\AgdaSpace{}%
\AgdaSymbol{:}\AgdaSpace{}%
\AgdaPrimitive{Set}\AgdaSpace{}%
\AgdaBound{i}\<%
\\
%
\>[6]\AgdaField{Sub}\AgdaSpace{}%
\AgdaSymbol{:}\AgdaSpace{}%
\AgdaField{Con}\AgdaSpace{}%
\AgdaSymbol{→}\AgdaSpace{}%
\AgdaField{Con}\AgdaSpace{}%
\AgdaSymbol{→}\AgdaSpace{}%
\AgdaPrimitive{Set}\AgdaSpace{}%
\AgdaSymbol{(}\AgdaBound{k}\AgdaSpace{}%
\AgdaOperator{\AgdaPrimitive{⊔}}\AgdaSpace{}%
\AgdaBound{l}\AgdaSymbol{)}\<%
\end{code}
A kategóriákban a morfizmusoknak kell rendelkezniük egy kompozíciós művelettel (\AgdaField{\_∘\_}) is, amely két morfizmus összekapcsolását végzi el, ezzel létrehozva egy új morfizmust. 
\begin{code}%
%
\>[6]\AgdaOperator{\AgdaField{\AgdaUnderscore{}∘\AgdaUnderscore{}}}\AgdaSpace{}%
\AgdaSymbol{:}\AgdaSpace{}%
\AgdaSymbol{∀\{}\AgdaBound{Γ}\AgdaSpace{}%
\AgdaBound{Δ}\AgdaSpace{}%
\AgdaBound{Θ}\AgdaSymbol{\}}\AgdaSpace{}%
\AgdaSymbol{→}\AgdaSpace{}%
\AgdaField{Sub}\AgdaSpace{}%
\AgdaBound{Δ}\AgdaSpace{}%
\AgdaBound{Γ}\AgdaSpace{}%
\AgdaSymbol{→}\AgdaSpace{}%
\AgdaField{Sub}\AgdaSpace{}%
\AgdaBound{Θ}\AgdaSpace{}%
\AgdaBound{Δ}\AgdaSpace{}%
\AgdaSymbol{→}\AgdaSpace{}%
\AgdaField{Sub}\AgdaSpace{}%
\AgdaBound{Θ}\AgdaSpace{}%
\AgdaBound{Γ}\<%
\\
%
\>[6]\AgdaField{ass}\AgdaSpace{}%
\AgdaSymbol{:}%
\>[58I]\AgdaSymbol{∀\{}\AgdaBound{Γ}\AgdaSpace{}%
\AgdaBound{Δ}\AgdaSpace{}%
\AgdaBound{Θ}\AgdaSpace{}%
\AgdaBound{Ξ}\AgdaSymbol{\}\{}\AgdaBound{γ}\AgdaSpace{}%
\AgdaSymbol{:}\AgdaSpace{}%
\AgdaField{Sub}\AgdaSpace{}%
\AgdaBound{Δ}\AgdaSpace{}%
\AgdaBound{Γ}\AgdaSymbol{\}\{}\AgdaBound{δ}\AgdaSpace{}%
\AgdaSymbol{:}\AgdaSpace{}%
\AgdaField{Sub}\AgdaSpace{}%
\AgdaBound{Θ}\AgdaSpace{}%
\AgdaBound{Δ}\AgdaSymbol{\}\{}\AgdaBound{θ}\AgdaSpace{}%
\AgdaSymbol{:}\AgdaSpace{}%
\AgdaField{Sub}\AgdaSpace{}%
\AgdaBound{Ξ}\AgdaSpace{}%
\AgdaBound{Θ}\AgdaSymbol{\}}\<%
\\
\>[.][@{}l@{}]\<[58I]%
\>[12]\AgdaSymbol{→}\AgdaSpace{}%
\AgdaSymbol{(}\AgdaBound{γ}\AgdaSpace{}%
\AgdaOperator{\AgdaField{∘}}\AgdaSpace{}%
\AgdaBound{δ}\AgdaSymbol{)}\AgdaSpace{}%
\AgdaOperator{\AgdaField{∘}}\AgdaSpace{}%
\AgdaBound{θ}\AgdaSpace{}%
\AgdaOperator{\AgdaDatatype{≡}}\AgdaSpace{}%
\AgdaBound{γ}\AgdaSpace{}%
\AgdaOperator{\AgdaField{∘}}\AgdaSpace{}%
\AgdaSymbol{(}\AgdaBound{δ}\AgdaSpace{}%
\AgdaOperator{\AgdaField{∘}}\AgdaSpace{}%
\AgdaBound{θ}\AgdaSymbol{)}\<%
\end{code}
Tehát, ha van egy morfizmus a Δ és Γ kontextusok között, valamint egy másik morfizmus a Θ és Δ között, akkor a kompozíciójuk egy új morfizmust ad a Θ és Γ kontextusok között. A kompozíciónak továbbá asszociatívnak kell lennie, amit az \AgdaField{ass} fejez ki.

A kategóriában minden objektumhoz léteznie kell egy olyan morfizmusnak, amely az objektumot önmagába képezi, ez lesz az identitás. Az identitás morfizmusnak a kompozícióval való alkalmazása nem változtathatja meg a kompozícióban szereplő másik morfizmust. Ezt az \AgdaField{idl} és \AgdaField{idr} egyenlőség biztosítja.
\begin{code}%
%
\>[6]\AgdaField{id}%
\>[10]\AgdaSymbol{:}\AgdaSpace{}%
\AgdaSymbol{∀\{}\AgdaBound{Γ}\AgdaSymbol{\}}\AgdaSpace{}%
\AgdaSymbol{→}\AgdaSpace{}%
\AgdaField{Sub}\AgdaSpace{}%
\AgdaBound{Γ}\AgdaSpace{}%
\AgdaBound{Γ}\<%
\\
%
\>[6]\AgdaField{idl}\AgdaSpace{}%
\AgdaSymbol{:}\AgdaSpace{}%
\AgdaSymbol{∀\{}\AgdaBound{Γ}\AgdaSpace{}%
\AgdaBound{Δ}\AgdaSymbol{\}\{}\AgdaBound{γ}\AgdaSpace{}%
\AgdaSymbol{:}\AgdaSpace{}%
\AgdaField{Sub}\AgdaSpace{}%
\AgdaBound{Δ}\AgdaSpace{}%
\AgdaBound{Γ}\AgdaSymbol{\}}\AgdaSpace{}%
\AgdaSymbol{→}\AgdaSpace{}%
\AgdaField{id}\AgdaSpace{}%
\AgdaOperator{\AgdaField{∘}}\AgdaSpace{}%
\AgdaBound{γ}\AgdaSpace{}%
\AgdaOperator{\AgdaDatatype{≡}}\AgdaSpace{}%
\AgdaBound{γ}\<%
\\
%
\>[6]\AgdaField{idr}\AgdaSpace{}%
\AgdaSymbol{:}\AgdaSpace{}%
\AgdaSymbol{∀\{}\AgdaBound{Γ}\AgdaSpace{}%
\AgdaBound{Δ}\AgdaSymbol{\}\{}\AgdaBound{γ}\AgdaSpace{}%
\AgdaSymbol{:}\AgdaSpace{}%
\AgdaField{Sub}\AgdaSpace{}%
\AgdaBound{Δ}\AgdaSpace{}%
\AgdaBound{Γ}\AgdaSymbol{\}}\AgdaSpace{}%
\AgdaSymbol{→}\AgdaSpace{}%
\AgdaBound{γ}\AgdaSpace{}%
\AgdaOperator{\AgdaField{∘}}\AgdaSpace{}%
\AgdaField{id}\AgdaSpace{}%
\AgdaOperator{\AgdaDatatype{≡}}\AgdaSpace{}%
\AgdaBound{γ}\<%
\end{code}
A terminális objektum egy olyan objektum, amelyhez minden más objektumból létezik pontosan egy morfizmus. A mi esetünkben ez a terminális objektum a \AgdaField{◇}, a hozzá tartozó morfizmus pedig az \AgdaField{ε}. A morfizmus egyedülállóságát a \AgdaField{◇η} fogalmazza meg.
\begin{code}%
%
\>[6]\AgdaField{◇}%
\>[10]\AgdaSymbol{:}\AgdaSpace{}%
\AgdaField{Con}\<%
\\
%
\>[6]\AgdaField{ε}%
\>[10]\AgdaSymbol{:}\AgdaSpace{}%
\AgdaSymbol{∀\{}\AgdaBound{Γ}\AgdaSymbol{\}}\AgdaSpace{}%
\AgdaSymbol{→}\AgdaSpace{}%
\AgdaField{Sub}\AgdaSpace{}%
\AgdaBound{Γ}\AgdaSpace{}%
\AgdaField{◇}\<%
\\
%
\>[6]\AgdaField{◇η}%
\>[10]\AgdaSymbol{:}\AgdaSpace{}%
\AgdaSymbol{∀\{}\AgdaBound{Γ}\AgdaSymbol{\}\{}\AgdaBound{σ}\AgdaSpace{}%
\AgdaSymbol{:}\AgdaSpace{}%
\AgdaField{Sub}\AgdaSpace{}%
\AgdaBound{Γ}\AgdaSpace{}%
\AgdaField{◇}\AgdaSymbol{\}}\AgdaSpace{}%
\AgdaSymbol{→}\AgdaSpace{}%
\AgdaBound{σ}\AgdaSpace{}%
\AgdaOperator{\AgdaDatatype{≡}}\AgdaSpace{}%
\AgdaField{ε}\<%
\end{code}

\section{Termek}

A termeket (\AgdaField{Tm}) egy funktorként adjuk meg. Úgy is mondhatjuk, hogy egy kontextustól függenek, mivel lehetnek benne termváltozók, amikbe be tudunk majd helyettesíteni, így értéket adva a szabad változóknak. Ehhez szükségünk lesz termekbe való behelyettesítésre, másnéven egy morfizmusok feletti akcióra. Egy funktornak teljesítenie kell két feltételt is: a kompozícióval való kompatibilitást és az identitással való kompatibilitást. Ezeket az \AgdaField{[∘]ᵗ} és \AgdaField{[id]ᵗ} egyenlőségek biztosítják.
\begin{code}%
%
\>[6]\AgdaField{Tm}%
\>[12]\AgdaSymbol{:}\AgdaSpace{}%
\AgdaField{Con}\AgdaSpace{}%
\AgdaSymbol{→}\AgdaSpace{}%
\AgdaPrimitive{Set}\AgdaSpace{}%
\AgdaBound{k}\<%
\\
%
\>[6]\AgdaOperator{\AgdaField{\AgdaUnderscore{}[\AgdaUnderscore{}]ᵗ}}\AgdaSpace{}%
\AgdaSymbol{:}\AgdaSpace{}%
\AgdaSymbol{∀\{}\AgdaBound{Γ}\AgdaSpace{}%
\AgdaBound{Δ}\AgdaSymbol{\}}\AgdaSpace{}%
\AgdaSymbol{→}\AgdaSpace{}%
\AgdaField{Tm}\AgdaSpace{}%
\AgdaBound{Γ}\AgdaSpace{}%
\AgdaSymbol{→}\AgdaSpace{}%
\AgdaField{Sub}\AgdaSpace{}%
\AgdaBound{Δ}\AgdaSpace{}%
\AgdaBound{Γ}\AgdaSpace{}%
\AgdaSymbol{→}\AgdaSpace{}%
\AgdaField{Tm}\AgdaSpace{}%
\AgdaBound{Δ}\<%
\\
%
\>[6]\AgdaField{[∘]ᵗ}%
\>[12]\AgdaSymbol{:}%
\>[148I]\AgdaSymbol{∀\{}\AgdaBound{Γ}\AgdaSpace{}%
\AgdaBound{Δ}\AgdaSpace{}%
\AgdaBound{θ}\AgdaSymbol{\}\{}\AgdaBound{t}\AgdaSpace{}%
\AgdaSymbol{:}\AgdaSpace{}%
\AgdaField{Tm}\AgdaSpace{}%
\AgdaBound{Γ}\AgdaSymbol{\}\{}\AgdaBound{γ}\AgdaSpace{}%
\AgdaSymbol{:}\AgdaSpace{}%
\AgdaField{Sub}\AgdaSpace{}%
\AgdaBound{Δ}\AgdaSpace{}%
\AgdaBound{Γ}\AgdaSymbol{\}\{}\AgdaBound{δ}\AgdaSpace{}%
\AgdaSymbol{:}\AgdaSpace{}%
\AgdaField{Sub}\AgdaSpace{}%
\AgdaBound{θ}\AgdaSpace{}%
\AgdaBound{Δ}\AgdaSymbol{\}}\<%
\\
\>[.][@{}l@{}]\<[148I]%
\>[14]\AgdaSymbol{→}\AgdaSpace{}%
\AgdaBound{t}\AgdaSpace{}%
\AgdaOperator{\AgdaField{[}}\AgdaSpace{}%
\AgdaBound{γ}\AgdaSpace{}%
\AgdaOperator{\AgdaField{∘}}\AgdaSpace{}%
\AgdaBound{δ}\AgdaSpace{}%
\AgdaOperator{\AgdaField{]ᵗ}}\AgdaSpace{}%
\AgdaOperator{\AgdaDatatype{≡}}\AgdaSpace{}%
\AgdaBound{t}\AgdaSpace{}%
\AgdaOperator{\AgdaField{[}}\AgdaSpace{}%
\AgdaBound{γ}\AgdaSpace{}%
\AgdaOperator{\AgdaField{]ᵗ}}\AgdaSpace{}%
\AgdaOperator{\AgdaField{[}}\AgdaSpace{}%
\AgdaBound{δ}\AgdaSpace{}%
\AgdaOperator{\AgdaField{]ᵗ}}\<%
\\
%
\>[6]\AgdaField{[id]ᵗ}\AgdaSpace{}%
\AgdaSymbol{:}\AgdaSpace{}%
\AgdaSymbol{∀\{}\AgdaBound{Γ}\AgdaSymbol{\}\{}\AgdaBound{t}\AgdaSpace{}%
\AgdaSymbol{:}\AgdaSpace{}%
\AgdaField{Tm}\AgdaSpace{}%
\AgdaBound{Γ}\AgdaSymbol{\}}\AgdaSpace{}%
\AgdaSymbol{→}\AgdaSpace{}%
\AgdaBound{t}\AgdaSpace{}%
\AgdaOperator{\AgdaField{[}}\AgdaSpace{}%
\AgdaField{id}\AgdaSpace{}%
\AgdaOperator{\AgdaField{]ᵗ}}\AgdaSpace{}%
\AgdaOperator{\AgdaDatatype{≡}}\AgdaSpace{}%
\AgdaBound{t}\<%
\end{code}
A környezetek termekkel és termváltozókkal való kiegészítését is megadjuk. A környezetbe betenni egy termváltozót (\AgdaField{\_▹ₜ}) hasonlítható egy olyan Descartes szorzathoz, amelynek az egyik oldalán az egyelemű halmaz található: \AgdaDatatype{\_×⊤}. A behelyettesítéseket is kiegészíthetjük termekkel (\AgdaField{\_,ₜ\_}), amelyeket ismételten a Descartes szorzattal lehetne azonosítani. A képzett pároknak tehát meg tudjuk adni az első (\AgdaField{pₜ}) és a második (\AgdaField{qₜ}) projekcióját.
\begin{code}%
%
\>[6]\AgdaOperator{\AgdaField{\AgdaUnderscore{}▹ₜ}}%
\>[12]\AgdaSymbol{:}\AgdaSpace{}%
\AgdaField{Con}\AgdaSpace{}%
\AgdaSymbol{→}\AgdaSpace{}%
\AgdaField{Con}\<%
\\
%
\>[6]\AgdaOperator{\AgdaField{\AgdaUnderscore{},ₜ\AgdaUnderscore{}}}%
\>[12]\AgdaSymbol{:}\AgdaSpace{}%
\AgdaSymbol{∀\{}\AgdaBound{Γ}\AgdaSpace{}%
\AgdaBound{Δ}\AgdaSymbol{\}}\AgdaSpace{}%
\AgdaSymbol{→}\AgdaSpace{}%
\AgdaField{Sub}\AgdaSpace{}%
\AgdaBound{Δ}\AgdaSpace{}%
\AgdaBound{Γ}\AgdaSpace{}%
\AgdaSymbol{→}\AgdaSpace{}%
\AgdaField{Tm}\AgdaSpace{}%
\AgdaBound{Δ}\AgdaSpace{}%
\AgdaSymbol{→}\AgdaSpace{}%
\AgdaField{Sub}\AgdaSpace{}%
\AgdaBound{Δ}\AgdaSpace{}%
\AgdaSymbol{(}\AgdaBound{Γ}\AgdaSpace{}%
\AgdaOperator{\AgdaField{▹ₜ}}\AgdaSymbol{)}\<%
\\
%
\>[6]\AgdaField{pₜ}%
\>[12]\AgdaSymbol{:}\AgdaSpace{}%
\AgdaSymbol{∀\{}\AgdaBound{Γ}\AgdaSymbol{\}}\AgdaSpace{}%
\AgdaSymbol{→}\AgdaSpace{}%
\AgdaField{Sub}\AgdaSpace{}%
\AgdaSymbol{(}\AgdaBound{Γ}\AgdaSpace{}%
\AgdaOperator{\AgdaField{▹ₜ}}\AgdaSymbol{)}\AgdaSpace{}%
\AgdaBound{Γ}\<%
\\
%
\>[6]\AgdaField{qₜ}%
\>[12]\AgdaSymbol{:}\AgdaSpace{}%
\AgdaSymbol{∀\{}\AgdaBound{Γ}\AgdaSymbol{\}}\AgdaSpace{}%
\AgdaSymbol{→}\AgdaSpace{}%
\AgdaField{Tm}\AgdaSpace{}%
\AgdaSymbol{(}\AgdaBound{Γ}\AgdaSpace{}%
\AgdaOperator{\AgdaField{▹ₜ}}\AgdaSymbol{)}\<%
\end{code}

Fontos megadjuk a redukciós szabályokat is (β, η), hogy biztosítsuk az azonos kifejezések közötti egyenlőséget.

\texttt{
((p∘\_) , (q[\_])) : Sub Δ (Γ ▹ₜ) ≅ Sub Δ Γ × Tm Δ : \_,ₜ\_
}

\begin{itemize}
\item β: jobbról balra, majd jobbra = mintha nem történt volna semmi
\item η: balról jobbra, majd balra = mintha nem történt volna semmi
\end{itemize}
\begin{code}%
%
\>[6]\AgdaField{▹ₜβ₁}%
\>[12]\AgdaSymbol{:}\AgdaSpace{}%
\AgdaSymbol{∀\{}\AgdaBound{Γ}\AgdaSpace{}%
\AgdaBound{Δ}\AgdaSymbol{\}\{}\AgdaBound{t}\AgdaSpace{}%
\AgdaSymbol{:}\AgdaSpace{}%
\AgdaField{Tm}\AgdaSpace{}%
\AgdaBound{Δ}\AgdaSymbol{\}\{}\AgdaBound{γ}\AgdaSpace{}%
\AgdaSymbol{:}\AgdaSpace{}%
\AgdaField{Sub}\AgdaSpace{}%
\AgdaBound{Δ}\AgdaSpace{}%
\AgdaBound{Γ}\AgdaSymbol{\}}\AgdaSpace{}%
\AgdaSymbol{→}\AgdaSpace{}%
\AgdaField{pₜ}\AgdaSpace{}%
\AgdaOperator{\AgdaField{∘}}\AgdaSpace{}%
\AgdaSymbol{(}\AgdaBound{γ}\AgdaSpace{}%
\AgdaOperator{\AgdaField{,ₜ}}\AgdaSpace{}%
\AgdaBound{t}\AgdaSymbol{)}\AgdaSpace{}%
\AgdaOperator{\AgdaDatatype{≡}}\AgdaSpace{}%
\AgdaBound{γ}\<%
\\
%
\>[6]\AgdaField{▹ₜβ₂}%
\>[12]\AgdaSymbol{:}\AgdaSpace{}%
\AgdaSymbol{∀\{}\AgdaBound{Γ}\AgdaSpace{}%
\AgdaBound{Δ}\AgdaSymbol{\}\{}\AgdaBound{t}\AgdaSpace{}%
\AgdaSymbol{:}\AgdaSpace{}%
\AgdaField{Tm}\AgdaSpace{}%
\AgdaBound{Δ}\AgdaSymbol{\}\{}\AgdaBound{γ}\AgdaSpace{}%
\AgdaSymbol{:}\AgdaSpace{}%
\AgdaField{Sub}\AgdaSpace{}%
\AgdaBound{Δ}\AgdaSpace{}%
\AgdaBound{Γ}\AgdaSymbol{\}}\AgdaSpace{}%
\AgdaSymbol{→}\AgdaSpace{}%
\AgdaField{qₜ}\AgdaSpace{}%
\AgdaOperator{\AgdaField{[}}\AgdaSpace{}%
\AgdaBound{γ}\AgdaSpace{}%
\AgdaOperator{\AgdaField{,ₜ}}\AgdaSpace{}%
\AgdaBound{t}\AgdaSpace{}%
\AgdaOperator{\AgdaField{]ᵗ}}\AgdaSpace{}%
\AgdaOperator{\AgdaDatatype{≡}}\AgdaSpace{}%
\AgdaBound{t}\<%
\\
%
\>[6]\AgdaField{▹ₜη}%
\>[12]\AgdaSymbol{:}\AgdaSpace{}%
\AgdaSymbol{∀\{}\AgdaBound{Γ}\AgdaSpace{}%
\AgdaBound{Δ}\AgdaSymbol{\}}\AgdaSpace{}%
\AgdaSymbol{→}\AgdaSpace{}%
\AgdaSymbol{\{}\AgdaBound{γt}\AgdaSpace{}%
\AgdaSymbol{:}\AgdaSpace{}%
\AgdaField{Sub}\AgdaSpace{}%
\AgdaBound{Δ}\AgdaSpace{}%
\AgdaSymbol{(}\AgdaBound{Γ}\AgdaSpace{}%
\AgdaOperator{\AgdaField{▹ₜ}}\AgdaSymbol{)\}}\AgdaSpace{}%
\AgdaSymbol{→}\AgdaSpace{}%
\AgdaBound{γt}\AgdaSpace{}%
\AgdaOperator{\AgdaDatatype{≡}}\AgdaSpace{}%
\AgdaSymbol{(}\AgdaField{pₜ}\AgdaSpace{}%
\AgdaOperator{\AgdaField{∘}}\AgdaSpace{}%
\AgdaBound{γt}\AgdaSpace{}%
\AgdaOperator{\AgdaField{,ₜ}}\AgdaSpace{}%
\AgdaField{qₜ}\AgdaSpace{}%
\AgdaOperator{\AgdaField{[}}\AgdaSpace{}%
\AgdaBound{γt}\AgdaSpace{}%
\AgdaOperator{\AgdaField{]ᵗ}}\AgdaSymbol{)}\<%
\end{code}

\section{Formulák}

A formulákra is tekinthetünk funktorként, így hasonlóan a termekhez a következő módon adjuk meg:
\begin{code}%
%
\>[6]\AgdaField{For}\AgdaSpace{}%
\AgdaSymbol{:}\AgdaSpace{}%
\AgdaField{Con}\AgdaSpace{}%
\AgdaSymbol{→}\AgdaSpace{}%
\AgdaPrimitive{Set}\AgdaSpace{}%
\AgdaBound{j}\<%
\\
%
\>[6]\AgdaOperator{\AgdaField{\AgdaUnderscore{}[\AgdaUnderscore{}]ᶠ}}\AgdaSpace{}%
\AgdaSymbol{:}\AgdaSpace{}%
\AgdaSymbol{∀\{}\AgdaBound{Γ}\AgdaSpace{}%
\AgdaBound{Δ}\AgdaSymbol{\}}\AgdaSpace{}%
\AgdaSymbol{→}\AgdaSpace{}%
\AgdaField{For}\AgdaSpace{}%
\AgdaBound{Γ}\AgdaSpace{}%
\AgdaSymbol{→}\AgdaSpace{}%
\AgdaField{Sub}\AgdaSpace{}%
\AgdaBound{Δ}\AgdaSpace{}%
\AgdaBound{Γ}\AgdaSpace{}%
\AgdaSymbol{→}\AgdaSpace{}%
\AgdaField{For}\AgdaSpace{}%
\AgdaBound{Δ}\<%
\\
%
\>[6]\AgdaField{[∘]ᶠ}%
\>[12]\AgdaSymbol{:}%
\>[289I]\AgdaSymbol{∀\{}\AgdaBound{Γ}\AgdaSpace{}%
\AgdaBound{Δ}\AgdaSpace{}%
\AgdaBound{θ}\AgdaSymbol{\}\{}\AgdaBound{A}\AgdaSpace{}%
\AgdaSymbol{:}\AgdaSpace{}%
\AgdaField{For}\AgdaSpace{}%
\AgdaBound{Γ}\AgdaSymbol{\}\{}\AgdaBound{γ}\AgdaSpace{}%
\AgdaSymbol{:}\AgdaSpace{}%
\AgdaField{Sub}\AgdaSpace{}%
\AgdaBound{Δ}\AgdaSpace{}%
\AgdaBound{Γ}\AgdaSymbol{\}\{}\AgdaBound{δ}\AgdaSpace{}%
\AgdaSymbol{:}\AgdaSpace{}%
\AgdaField{Sub}\AgdaSpace{}%
\AgdaBound{θ}\AgdaSpace{}%
\AgdaBound{Δ}\AgdaSymbol{\}}\<%
\\
\>[.][@{}l@{}]\<[289I]%
\>[14]\AgdaSymbol{→}\AgdaSpace{}%
\AgdaBound{A}\AgdaSpace{}%
\AgdaOperator{\AgdaField{[}}\AgdaSpace{}%
\AgdaBound{γ}\AgdaSpace{}%
\AgdaOperator{\AgdaField{∘}}\AgdaSpace{}%
\AgdaBound{δ}\AgdaSpace{}%
\AgdaOperator{\AgdaField{]ᶠ}}\AgdaSpace{}%
\AgdaOperator{\AgdaDatatype{≡}}\AgdaSpace{}%
\AgdaBound{A}\AgdaSpace{}%
\AgdaOperator{\AgdaField{[}}\AgdaSpace{}%
\AgdaBound{γ}\AgdaSpace{}%
\AgdaOperator{\AgdaField{]ᶠ}}\AgdaSpace{}%
\AgdaOperator{\AgdaField{[}}\AgdaSpace{}%
\AgdaBound{δ}\AgdaSpace{}%
\AgdaOperator{\AgdaField{]ᶠ}}\<%
\\
%
\>[6]\AgdaField{[id]ᶠ}\AgdaSpace{}%
\AgdaSymbol{:}\AgdaSpace{}%
\AgdaSymbol{∀\{}\AgdaBound{Γ}\AgdaSymbol{\}\{}\AgdaBound{A}\AgdaSpace{}%
\AgdaSymbol{:}\AgdaSpace{}%
\AgdaField{For}\AgdaSpace{}%
\AgdaBound{Γ}\AgdaSymbol{\}}\AgdaSpace{}%
\AgdaSymbol{→}\AgdaSpace{}%
\AgdaBound{A}\AgdaSpace{}%
\AgdaOperator{\AgdaField{[}}\AgdaSpace{}%
\AgdaField{id}\AgdaSpace{}%
\AgdaOperator{\AgdaField{]ᶠ}}\AgdaSpace{}%
\AgdaOperator{\AgdaDatatype{≡}}\AgdaSpace{}%
\AgdaBound{A}\<%
\end{code}

\section{Bizonyítások}

A bizonyításokat már egy kicsivel bonyolultabban adjuk meg, ugyanis egy \AgdaField{For} feletti dependáns funktronak is nevezhetnénk. Ha vesszük például a Γ ⊢ A szokásos logikai (bizonyításelméleti) jelölést, akkor a \AgdaDatatype{p} : \AgdaField{Pf} Γ A azt jelenti, hogy \AgdaDatatype{p} bizonyítja az A állítást, ahol a szabad változók Γ-ban vannak.
\begin{code}%
%
\>[6]\AgdaField{Pf}%
\>[10]\AgdaSymbol{:}\AgdaSpace{}%
\AgdaSymbol{(}\AgdaBound{Γ}\AgdaSpace{}%
\AgdaSymbol{:}\AgdaSpace{}%
\AgdaField{Con}\AgdaSymbol{)}\AgdaSpace{}%
\AgdaSymbol{→}\AgdaSpace{}%
\AgdaField{For}\AgdaSpace{}%
\AgdaBound{Γ}\AgdaSpace{}%
\AgdaSymbol{→}\AgdaSpace{}%
\AgdaPrimitive{Prop}\AgdaSpace{}%
\AgdaBound{l}\<%
\\
%
\>[6]\AgdaOperator{\AgdaField{\AgdaUnderscore{}[\AgdaUnderscore{}]ᵖ}}\AgdaSpace{}%
\AgdaSymbol{:}\AgdaSpace{}%
\AgdaSymbol{∀\{}\AgdaBound{Γ}\AgdaSpace{}%
\AgdaBound{Δ}\AgdaSpace{}%
\AgdaBound{A}\AgdaSymbol{\}}\AgdaSpace{}%
\AgdaSymbol{→}\AgdaSpace{}%
\AgdaField{Pf}\AgdaSpace{}%
\AgdaBound{Γ}\AgdaSpace{}%
\AgdaBound{A}\AgdaSpace{}%
\AgdaSymbol{→}\AgdaSpace{}%
\AgdaSymbol{(}\AgdaBound{γ}\AgdaSpace{}%
\AgdaSymbol{:}\AgdaSpace{}%
\AgdaField{Sub}\AgdaSpace{}%
\AgdaBound{Δ}\AgdaSpace{}%
\AgdaBound{Γ}\AgdaSymbol{)}\AgdaSpace{}%
\AgdaSymbol{→}\AgdaSpace{}%
\AgdaField{Pf}\AgdaSpace{}%
\AgdaBound{Δ}\AgdaSpace{}%
\AgdaSymbol{(}\AgdaBound{A}\AgdaSpace{}%
\AgdaOperator{\AgdaField{[}}\AgdaSpace{}%
\AgdaBound{γ}\AgdaSpace{}%
\AgdaOperator{\AgdaField{]ᶠ}}\AgdaSymbol{)}\<%
\end{code}
A bizonyítás-változókat a termváltozók esetében is használt operációkhoz hasonlóan valósítjuk meg:
\begin{code}%
%
\>[6]\AgdaOperator{\AgdaField{\AgdaUnderscore{}▹ₚ\AgdaUnderscore{}}}\AgdaSpace{}%
\AgdaSymbol{:}\AgdaSpace{}%
\AgdaSymbol{(}\AgdaBound{Γ}\AgdaSpace{}%
\AgdaSymbol{:}\AgdaSpace{}%
\AgdaField{Con}\AgdaSymbol{)}\AgdaSpace{}%
\AgdaSymbol{→}\AgdaSpace{}%
\AgdaField{For}\AgdaSpace{}%
\AgdaBound{Γ}\AgdaSpace{}%
\AgdaSymbol{→}\AgdaSpace{}%
\AgdaField{Con}\<%
\\
%
\>[6]\AgdaOperator{\AgdaField{\AgdaUnderscore{},ₚ\AgdaUnderscore{}}}%
\>[12]\AgdaSymbol{:}\AgdaSpace{}%
\AgdaSymbol{∀\{}\AgdaBound{Δ}\AgdaSpace{}%
\AgdaBound{Γ}\AgdaSpace{}%
\AgdaBound{A}\AgdaSymbol{\}}\AgdaSpace{}%
\AgdaSymbol{→}\AgdaSpace{}%
\AgdaSymbol{(}\AgdaBound{γ}\AgdaSpace{}%
\AgdaSymbol{:}\AgdaSpace{}%
\AgdaField{Sub}\AgdaSpace{}%
\AgdaBound{Δ}\AgdaSpace{}%
\AgdaBound{Γ}\AgdaSymbol{)}\AgdaSpace{}%
\AgdaSymbol{→}\AgdaSpace{}%
\AgdaField{Pf}\AgdaSpace{}%
\AgdaBound{Δ}\AgdaSpace{}%
\AgdaSymbol{(}\AgdaBound{A}\AgdaSpace{}%
\AgdaOperator{\AgdaField{[}}\AgdaSpace{}%
\AgdaBound{γ}\AgdaSpace{}%
\AgdaOperator{\AgdaField{]ᶠ}}\AgdaSymbol{)}\AgdaSpace{}%
\AgdaSymbol{→}\AgdaSpace{}%
\AgdaField{Sub}\AgdaSpace{}%
\AgdaBound{Δ}\AgdaSpace{}%
\AgdaSymbol{(}\AgdaBound{Γ}\AgdaSpace{}%
\AgdaOperator{\AgdaField{▹ₚ}}\AgdaSpace{}%
\AgdaBound{A}\AgdaSymbol{)}\<%
\\
%
\>[6]\AgdaField{pₚ}%
\>[12]\AgdaSymbol{:}\AgdaSpace{}%
\AgdaSymbol{∀\{}\AgdaBound{Γ}\AgdaSpace{}%
\AgdaBound{A}\AgdaSymbol{\}}\AgdaSpace{}%
\AgdaSymbol{→}\AgdaSpace{}%
\AgdaField{Sub}\AgdaSpace{}%
\AgdaSymbol{(}\AgdaBound{Γ}\AgdaSpace{}%
\AgdaOperator{\AgdaField{▹ₚ}}\AgdaSpace{}%
\AgdaBound{A}\AgdaSymbol{)}\AgdaSpace{}%
\AgdaBound{Γ}\<%
\\
%
\>[6]\AgdaField{qₚ}%
\>[12]\AgdaSymbol{:}\AgdaSpace{}%
\AgdaSymbol{∀\{}\AgdaBound{Γ}\AgdaSpace{}%
\AgdaBound{A}\AgdaSymbol{\}}\AgdaSpace{}%
\AgdaSymbol{→}\AgdaSpace{}%
\AgdaField{Pf}\AgdaSpace{}%
\AgdaSymbol{(}\AgdaBound{Γ}\AgdaSpace{}%
\AgdaOperator{\AgdaField{▹ₚ}}\AgdaSpace{}%
\AgdaBound{A}\AgdaSymbol{)}\AgdaSpace{}%
\AgdaSymbol{((}\AgdaBound{A}\AgdaSpace{}%
\AgdaOperator{\AgdaField{[}}\AgdaSpace{}%
\AgdaField{pₚ}\AgdaSpace{}%
\AgdaOperator{\AgdaField{]ᶠ}}\AgdaSymbol{))}\<%
\\
%
\>[6]\AgdaField{▹ₚβ₁}\AgdaSpace{}%
\AgdaSymbol{:}\AgdaSpace{}%
\AgdaSymbol{∀}\AgdaSpace{}%
\AgdaSymbol{\{}\AgdaBound{Γ}\AgdaSpace{}%
\AgdaBound{Δ}\AgdaSpace{}%
\AgdaBound{A}\AgdaSymbol{\}\{}\AgdaBound{γ}\AgdaSpace{}%
\AgdaSymbol{:}\AgdaSpace{}%
\AgdaField{Sub}\AgdaSpace{}%
\AgdaBound{Δ}\AgdaSpace{}%
\AgdaBound{Γ}\AgdaSymbol{\}\{}\AgdaBound{a}\AgdaSpace{}%
\AgdaSymbol{:}\AgdaSpace{}%
\AgdaField{Pf}\AgdaSpace{}%
\AgdaBound{Δ}\AgdaSpace{}%
\AgdaSymbol{(}\AgdaBound{A}\AgdaSpace{}%
\AgdaOperator{\AgdaField{[}}\AgdaSpace{}%
\AgdaBound{γ}\AgdaSpace{}%
\AgdaOperator{\AgdaField{]ᶠ}}\AgdaSymbol{)\}}\AgdaSpace{}%
\AgdaSymbol{→}\AgdaSpace{}%
\AgdaField{pₚ}\AgdaSpace{}%
\AgdaOperator{\AgdaField{∘}}\AgdaSpace{}%
\AgdaSymbol{(}\AgdaBound{γ}\AgdaSpace{}%
\AgdaOperator{\AgdaField{,ₚ}}\AgdaSpace{}%
\AgdaBound{a}\AgdaSymbol{)}\AgdaSpace{}%
\AgdaOperator{\AgdaDatatype{≡}}\AgdaSpace{}%
\AgdaBound{γ}\<%
\\
%
\>[6]\AgdaField{▹ₚη}%
\>[11]\AgdaSymbol{:}%
\>[433I]\AgdaSymbol{∀}\AgdaSpace{}%
\AgdaSymbol{\{}\AgdaBound{Γ}\AgdaSpace{}%
\AgdaBound{Δ}\AgdaSpace{}%
\AgdaBound{A}\AgdaSymbol{\}\{}\AgdaBound{γa}\AgdaSpace{}%
\AgdaSymbol{:}\AgdaSpace{}%
\AgdaField{Sub}\AgdaSpace{}%
\AgdaBound{Δ}\AgdaSpace{}%
\AgdaSymbol{(}\AgdaBound{Γ}\AgdaSpace{}%
\AgdaOperator{\AgdaField{▹ₚ}}\AgdaSpace{}%
\AgdaBound{A}\AgdaSymbol{)\}}\<%
\\
\>[.][@{}l@{}]\<[433I]%
\>[13]\AgdaSymbol{→}\AgdaSpace{}%
\AgdaBound{γa}\AgdaSpace{}%
\AgdaOperator{\AgdaDatatype{≡}}\AgdaSpace{}%
\AgdaSymbol{(}\AgdaField{pₚ}\AgdaSpace{}%
\AgdaOperator{\AgdaField{∘}}\AgdaSpace{}%
\AgdaBound{γa}\AgdaSpace{}%
\AgdaOperator{\AgdaField{,ₚ}}\AgdaSpace{}%
\AgdaFunction{substP}\AgdaSpace{}%
\AgdaSymbol{(}\AgdaField{Pf}\AgdaSpace{}%
\AgdaBound{Δ}\AgdaSymbol{)}\AgdaSpace{}%
\AgdaSymbol{(}\AgdaField{[∘]ᶠ}\AgdaSpace{}%
\AgdaOperator{\AgdaFunction{⁻¹}}\AgdaSymbol{)}\AgdaSpace{}%
\AgdaSymbol{(}\AgdaField{qₚ}\AgdaSpace{}%
\AgdaOperator{\AgdaField{[}}\AgdaSpace{}%
\AgdaBound{γa}\AgdaSpace{}%
\AgdaOperator{\AgdaField{]ᵖ}}\AgdaSymbol{))}\<%
\end{code}

\noindent
\raggedright
\texttt{
((p∘\_),(q[\_])) : Sub Δ (Γ▹ₚA) ≅ (γ: Sub Δ Γ) × Pf Δ (A[γ]ᶠ) : \_,ₚ\_
}

\section{Reláció- és függvényszimbólumok}

Reláció és függvényszimbólumok (a modell a függvények és a relációk aritásával van felparaméterezve):
\begin{code}%
%
\>[6]\AgdaField{Rel}%
\>[12]\AgdaSymbol{:}\AgdaSpace{}%
\AgdaSymbol{∀\{}\AgdaBound{Γ}\AgdaSymbol{\}\{}\AgdaBound{n}\AgdaSpace{}%
\AgdaSymbol{:}\AgdaSpace{}%
\AgdaDatatype{ℕ}\AgdaSymbol{\}}\AgdaSpace{}%
\AgdaSymbol{→}\AgdaSpace{}%
\AgdaBound{relar}\AgdaSpace{}%
\AgdaBound{n}\AgdaSpace{}%
\AgdaSymbol{→}\AgdaSpace{}%
\AgdaField{Tm}\AgdaSpace{}%
\AgdaBound{Γ}\AgdaSpace{}%
\AgdaOperator{\AgdaFunction{\textasciicircum{}}}\AgdaSpace{}%
\AgdaBound{n}\AgdaSpace{}%
\AgdaSymbol{→}\AgdaSpace{}%
\AgdaField{For}\AgdaSpace{}%
\AgdaBound{Γ}\<%
\\
%
\>[6]\AgdaField{Rel[]}\AgdaSpace{}%
\AgdaSymbol{:}%
\>[473I]\AgdaSymbol{∀\{}\AgdaBound{Γ}\AgdaSpace{}%
\AgdaBound{n}\AgdaSymbol{\}\{}\AgdaBound{ar}\AgdaSpace{}%
\AgdaSymbol{:}\AgdaSpace{}%
\AgdaBound{relar}\AgdaSpace{}%
\AgdaBound{n}\AgdaSymbol{\}\{}\AgdaBound{ts}\AgdaSpace{}%
\AgdaSymbol{:}\AgdaSpace{}%
\AgdaField{Tm}\AgdaSpace{}%
\AgdaBound{Γ}\AgdaSpace{}%
\AgdaOperator{\AgdaFunction{\textasciicircum{}}}\AgdaSpace{}%
\AgdaBound{n}\AgdaSymbol{\}\{}\AgdaBound{Δ}\AgdaSymbol{\}\{}\AgdaBound{γ}\AgdaSpace{}%
\AgdaSymbol{:}\AgdaSpace{}%
\AgdaField{Sub}\AgdaSpace{}%
\AgdaBound{Δ}\AgdaSpace{}%
\AgdaBound{Γ}\AgdaSymbol{\}}\<%
\\
\>[.][@{}l@{}]\<[473I]%
\>[14]\AgdaSymbol{→}\AgdaSpace{}%
\AgdaField{Rel}\AgdaSpace{}%
\AgdaBound{ar}\AgdaSpace{}%
\AgdaBound{ts}\AgdaSpace{}%
\AgdaOperator{\AgdaField{[}}\AgdaSpace{}%
\AgdaBound{γ}\AgdaSpace{}%
\AgdaOperator{\AgdaField{]ᶠ}}\AgdaSpace{}%
\AgdaOperator{\AgdaDatatype{≡}}\AgdaSpace{}%
\AgdaField{Rel}\AgdaSpace{}%
\AgdaBound{ar}\AgdaSpace{}%
\AgdaSymbol{(}\AgdaFunction{map}\AgdaSpace{}%
\AgdaOperator{\AgdaField{\AgdaUnderscore{}[}}\AgdaSpace{}%
\AgdaBound{γ}\AgdaSpace{}%
\AgdaOperator{\AgdaField{]ᵗ}}\AgdaSpace{}%
\AgdaBound{ts}\AgdaSymbol{)}\<%
\\
%
\>[6]\AgdaField{fun}%
\>[12]\AgdaSymbol{:}\AgdaSpace{}%
\AgdaSymbol{∀\{}\AgdaBound{Γ}\AgdaSymbol{\}\{}\AgdaBound{n}\AgdaSpace{}%
\AgdaSymbol{:}\AgdaSpace{}%
\AgdaDatatype{ℕ}\AgdaSymbol{\}}\AgdaSpace{}%
\AgdaSymbol{→}\AgdaSpace{}%
\AgdaBound{funar}\AgdaSpace{}%
\AgdaBound{n}\AgdaSpace{}%
\AgdaSymbol{→}\AgdaSpace{}%
\AgdaField{Tm}\AgdaSpace{}%
\AgdaBound{Γ}\AgdaSpace{}%
\AgdaOperator{\AgdaFunction{\textasciicircum{}}}\AgdaSpace{}%
\AgdaBound{n}\AgdaSpace{}%
\AgdaSymbol{→}\AgdaSpace{}%
\AgdaField{Tm}\AgdaSpace{}%
\AgdaBound{Γ}\<%
\\
%
\>[6]\AgdaField{fun[]}\AgdaSpace{}%
\AgdaSymbol{:}%
\>[516I]\AgdaSymbol{∀\{}\AgdaBound{Γ}\AgdaSpace{}%
\AgdaBound{n}\AgdaSymbol{\}\{}\AgdaBound{ar}\AgdaSpace{}%
\AgdaSymbol{:}\AgdaSpace{}%
\AgdaBound{funar}\AgdaSpace{}%
\AgdaBound{n}\AgdaSymbol{\}\{}\AgdaBound{ts}\AgdaSpace{}%
\AgdaSymbol{:}\AgdaSpace{}%
\AgdaField{Tm}\AgdaSpace{}%
\AgdaBound{Γ}\AgdaSpace{}%
\AgdaOperator{\AgdaFunction{\textasciicircum{}}}\AgdaSpace{}%
\AgdaBound{n}\AgdaSymbol{\}\{}\AgdaBound{Δ}\AgdaSymbol{\}\{}\AgdaBound{γ}\AgdaSpace{}%
\AgdaSymbol{:}\AgdaSpace{}%
\AgdaField{Sub}\AgdaSpace{}%
\AgdaBound{Δ}\AgdaSpace{}%
\AgdaBound{Γ}\AgdaSymbol{\}}\<%
\\
\>[.][@{}l@{}]\<[516I]%
\>[14]\AgdaSymbol{→}\AgdaSpace{}%
\AgdaField{fun}\AgdaSpace{}%
\AgdaBound{ar}\AgdaSpace{}%
\AgdaBound{ts}\AgdaSpace{}%
\AgdaOperator{\AgdaField{[}}\AgdaSpace{}%
\AgdaBound{γ}\AgdaSpace{}%
\AgdaOperator{\AgdaField{]ᵗ}}\AgdaSpace{}%
\AgdaOperator{\AgdaDatatype{≡}}\AgdaSpace{}%
\AgdaField{fun}\AgdaSpace{}%
\AgdaBound{ar}\AgdaSpace{}%
\AgdaSymbol{(}\AgdaFunction{map}\AgdaSpace{}%
\AgdaOperator{\AgdaField{\AgdaUnderscore{}[}}\AgdaSpace{}%
\AgdaBound{γ}\AgdaSpace{}%
\AgdaOperator{\AgdaField{]ᵗ}}\AgdaSpace{}%
\AgdaBound{ts}\AgdaSymbol{)}\<%
\end{code}

\section{Logikai összekötők}

\subsection{Implikáció (⊃)}

Az A ⊃ B jelentése: "ha A igaz, akkor B is igaz".
\begin{code}%
%
\>[6]\AgdaOperator{\AgdaField{\AgdaUnderscore{}⊃\AgdaUnderscore{}}}%
\>[12]\AgdaSymbol{:}\AgdaSpace{}%
\AgdaSymbol{∀\{}\AgdaBound{Γ}\AgdaSymbol{\}}\AgdaSpace{}%
\AgdaSymbol{→}\AgdaSpace{}%
\AgdaField{For}\AgdaSpace{}%
\AgdaBound{Γ}\AgdaSpace{}%
\AgdaSymbol{→}\AgdaSpace{}%
\AgdaField{For}\AgdaSpace{}%
\AgdaBound{Γ}\AgdaSpace{}%
\AgdaSymbol{→}\AgdaSpace{}%
\AgdaField{For}\AgdaSpace{}%
\AgdaBound{Γ}\<%
\end{code}
\begin{code}%
%
\>[6]\AgdaField{⊃[]}%
\>[12]\AgdaSymbol{:}\AgdaSpace{}%
\AgdaSymbol{∀\{}\AgdaBound{Γ}\AgdaSpace{}%
\AgdaBound{A}\AgdaSpace{}%
\AgdaBound{B}\AgdaSpace{}%
\AgdaBound{Δ}\AgdaSymbol{\}\{}\AgdaBound{γ}\AgdaSpace{}%
\AgdaSymbol{:}\AgdaSpace{}%
\AgdaField{Sub}\AgdaSpace{}%
\AgdaBound{Δ}\AgdaSpace{}%
\AgdaBound{Γ}\AgdaSymbol{\}}\AgdaSpace{}%
\AgdaSymbol{→}\AgdaSpace{}%
\AgdaSymbol{(}\AgdaBound{A}\AgdaSpace{}%
\AgdaOperator{\AgdaField{⊃}}\AgdaSpace{}%
\AgdaBound{B}\AgdaSymbol{)}\AgdaSpace{}%
\AgdaOperator{\AgdaField{[}}\AgdaSpace{}%
\AgdaBound{γ}\AgdaSpace{}%
\AgdaOperator{\AgdaField{]ᶠ}}\AgdaSpace{}%
\AgdaOperator{\AgdaDatatype{≡}}\AgdaSpace{}%
\AgdaBound{A}\AgdaSpace{}%
\AgdaOperator{\AgdaField{[}}\AgdaSpace{}%
\AgdaBound{γ}\AgdaSpace{}%
\AgdaOperator{\AgdaField{]ᶠ}}\AgdaSpace{}%
\AgdaOperator{\AgdaField{⊃}}\AgdaSpace{}%
\AgdaBound{B}\AgdaSpace{}%
\AgdaOperator{\AgdaField{[}}\AgdaSpace{}%
\AgdaBound{γ}\AgdaSpace{}%
\AgdaOperator{\AgdaField{]ᶠ}}\<%
\end{code}
\begin{itemize}
\item Az implikáció bevezető szabálya (\AgdaField{⊃in}): Ha tudjuk bizonyítani, hogy B igaz egy olyan kontextusban, ahol A feltételezett, akkor van egy bizonyításunk A ⊃ B-re is.
\begin{code}%
%
\>[6]\AgdaField{⊃in}%
\>[12]\AgdaSymbol{:}\AgdaSpace{}%
\AgdaSymbol{∀\{}\AgdaBound{Γ}\AgdaSpace{}%
\AgdaBound{A}\AgdaSpace{}%
\AgdaBound{B}\AgdaSymbol{\}}\AgdaSpace{}%
\AgdaSymbol{→}\AgdaSpace{}%
\AgdaField{Pf}\AgdaSpace{}%
\AgdaSymbol{(}\AgdaBound{Γ}\AgdaSpace{}%
\AgdaOperator{\AgdaField{▹ₚ}}\AgdaSpace{}%
\AgdaBound{A}\AgdaSymbol{)}\AgdaSpace{}%
\AgdaSymbol{(}\AgdaBound{B}\AgdaSpace{}%
\AgdaOperator{\AgdaField{[}}\AgdaSpace{}%
\AgdaField{pₚ}\AgdaSpace{}%
\AgdaOperator{\AgdaField{]ᶠ}}\AgdaSymbol{)}\AgdaSpace{}%
\AgdaSymbol{→}\AgdaSpace{}%
\AgdaField{Pf}\AgdaSpace{}%
\AgdaBound{Γ}\AgdaSpace{}%
\AgdaSymbol{((}\AgdaBound{A}\AgdaSpace{}%
\AgdaOperator{\AgdaField{⊃}}\AgdaSpace{}%
\AgdaBound{B}\AgdaSymbol{))}\<%
\end{code}
\item Az implikáció kivezető szabálya (\AgdaField{⊃out}): Ha van egy bizonyításunk A ⊃ B-re, valamint A-ra is, akkor van egy bizonyításunk B-re is.
\begin{code}%
%
\>[6]\AgdaField{⊃out}%
\>[12]\AgdaSymbol{:}\AgdaSpace{}%
\AgdaSymbol{∀\{}\AgdaBound{Γ}\AgdaSpace{}%
\AgdaBound{A}\AgdaSpace{}%
\AgdaBound{B}\AgdaSymbol{\}}\AgdaSpace{}%
\AgdaSymbol{→}\AgdaSpace{}%
\AgdaField{Pf}\AgdaSpace{}%
\AgdaBound{Γ}\AgdaSpace{}%
\AgdaSymbol{(}\AgdaBound{A}\AgdaSpace{}%
\AgdaOperator{\AgdaField{⊃}}\AgdaSpace{}%
\AgdaBound{B}\AgdaSymbol{)}\AgdaSpace{}%
\AgdaSymbol{→}\AgdaSpace{}%
\AgdaField{Pf}\AgdaSpace{}%
\AgdaBound{Γ}\AgdaSpace{}%
\AgdaBound{A}\AgdaSpace{}%
\AgdaSymbol{→}\AgdaSpace{}%
\AgdaField{Pf}\AgdaSpace{}%
\AgdaBound{Γ}\AgdaSpace{}%
\AgdaBound{B}\<%
\end{code}
\end{itemize}

\subsection{Konjunkció (∧)}
Az A ∧ B azt jelenti, hogy mind A, mind B igaz.
\begin{code}%
%
\>[6]\AgdaOperator{\AgdaField{\AgdaUnderscore{}∧\AgdaUnderscore{}}}%
\>[12]\AgdaSymbol{:}\AgdaSpace{}%
\AgdaSymbol{∀\{}\AgdaBound{Γ}\AgdaSymbol{\}}\AgdaSpace{}%
\AgdaSymbol{→}\AgdaSpace{}%
\AgdaField{For}\AgdaSpace{}%
\AgdaBound{Γ}\AgdaSpace{}%
\AgdaSymbol{→}\AgdaSpace{}%
\AgdaField{For}\AgdaSpace{}%
\AgdaBound{Γ}\AgdaSpace{}%
\AgdaSymbol{→}\AgdaSpace{}%
\AgdaField{For}\AgdaSpace{}%
\AgdaBound{Γ}\<%
\end{code}
\begin{code}%
%
\>[6]\AgdaField{∧[]}%
\>[12]\AgdaSymbol{:}\AgdaSpace{}%
\AgdaSymbol{∀\{}\AgdaBound{Γ}\AgdaSpace{}%
\AgdaBound{A}\AgdaSpace{}%
\AgdaBound{B}\AgdaSpace{}%
\AgdaBound{Δ}\AgdaSymbol{\}\{}\AgdaBound{γ}\AgdaSpace{}%
\AgdaSymbol{:}\AgdaSpace{}%
\AgdaField{Sub}\AgdaSpace{}%
\AgdaBound{Δ}\AgdaSpace{}%
\AgdaBound{Γ}\AgdaSymbol{\}}\AgdaSpace{}%
\AgdaSymbol{→}\AgdaSpace{}%
\AgdaSymbol{(}\AgdaBound{A}\AgdaSpace{}%
\AgdaOperator{\AgdaField{∧}}\AgdaSpace{}%
\AgdaBound{B}\AgdaSymbol{)}\AgdaSpace{}%
\AgdaOperator{\AgdaField{[}}\AgdaSpace{}%
\AgdaBound{γ}\AgdaSpace{}%
\AgdaOperator{\AgdaField{]ᶠ}}\AgdaSpace{}%
\AgdaOperator{\AgdaDatatype{≡}}\AgdaSpace{}%
\AgdaBound{A}\AgdaSpace{}%
\AgdaOperator{\AgdaField{[}}\AgdaSpace{}%
\AgdaBound{γ}\AgdaSpace{}%
\AgdaOperator{\AgdaField{]ᶠ}}\AgdaSpace{}%
\AgdaOperator{\AgdaField{∧}}\AgdaSpace{}%
\AgdaBound{B}\AgdaSpace{}%
\AgdaOperator{\AgdaField{[}}\AgdaSpace{}%
\AgdaBound{γ}\AgdaSpace{}%
\AgdaOperator{\AgdaField{]ᶠ}}\<%
\end{code}
\begin{itemize}
\item A konjunkció bevezető szabálya (\AgdaField{∧in}): Ha A és B külön-külön bizonyítható, akkor A ∧ B is igaz.
\begin{code}%
%
\>[6]\AgdaField{∧in}%
\>[12]\AgdaSymbol{:}\AgdaSpace{}%
\AgdaSymbol{∀\{}\AgdaBound{Γ}\AgdaSpace{}%
\AgdaBound{A}\AgdaSpace{}%
\AgdaBound{B}\AgdaSymbol{\}}\AgdaSpace{}%
\AgdaSymbol{→}\AgdaSpace{}%
\AgdaField{Pf}\AgdaSpace{}%
\AgdaBound{Γ}\AgdaSpace{}%
\AgdaBound{A}\AgdaSpace{}%
\AgdaSymbol{→}\AgdaSpace{}%
\AgdaField{Pf}\AgdaSpace{}%
\AgdaBound{Γ}\AgdaSpace{}%
\AgdaBound{B}\AgdaSpace{}%
\AgdaSymbol{→}\AgdaSpace{}%
\AgdaField{Pf}\AgdaSpace{}%
\AgdaBound{Γ}\AgdaSpace{}%
\AgdaSymbol{(}\AgdaBound{A}\AgdaSpace{}%
\AgdaOperator{\AgdaField{∧}}\AgdaSpace{}%
\AgdaBound{B}\AgdaSymbol{)}\<%
\end{code}
\item Az konjunkció kivezető szabálya (\AgdaField{∧out₁}): Az A ∧ B-ből következik A.
\item Az konjunkció kivezető szabálya (\AgdaField{∧out₂}): Az A ∧ B-ből következik B.
\begin{code}%
%
\>[6]\AgdaField{∧out₁}\AgdaSpace{}%
\AgdaSymbol{:}\AgdaSpace{}%
\AgdaSymbol{∀\{}\AgdaBound{Γ}\AgdaSpace{}%
\AgdaBound{A}\AgdaSpace{}%
\AgdaBound{B}\AgdaSymbol{\}}\AgdaSpace{}%
\AgdaSymbol{→}\AgdaSpace{}%
\AgdaField{Pf}\AgdaSpace{}%
\AgdaBound{Γ}\AgdaSpace{}%
\AgdaSymbol{(}\AgdaBound{A}\AgdaSpace{}%
\AgdaOperator{\AgdaField{∧}}\AgdaSpace{}%
\AgdaBound{B}\AgdaSymbol{)}\AgdaSpace{}%
\AgdaSymbol{→}\AgdaSpace{}%
\AgdaField{Pf}\AgdaSpace{}%
\AgdaBound{Γ}\AgdaSpace{}%
\AgdaBound{A}\<%
\\
%
\>[6]\AgdaField{∧out₂}\AgdaSpace{}%
\AgdaSymbol{:}\AgdaSpace{}%
\AgdaSymbol{∀\{}\AgdaBound{Γ}\AgdaSpace{}%
\AgdaBound{A}\AgdaSpace{}%
\AgdaBound{B}\AgdaSymbol{\}}\AgdaSpace{}%
\AgdaSymbol{→}\AgdaSpace{}%
\AgdaField{Pf}\AgdaSpace{}%
\AgdaBound{Γ}\AgdaSpace{}%
\AgdaSymbol{(}\AgdaBound{A}\AgdaSpace{}%
\AgdaOperator{\AgdaField{∧}}\AgdaSpace{}%
\AgdaBound{B}\AgdaSymbol{)}\AgdaSpace{}%
\AgdaSymbol{→}\AgdaSpace{}%
\AgdaField{Pf}\AgdaSpace{}%
\AgdaBound{Γ}\AgdaSpace{}%
\AgdaBound{B}\<%
\end{code}
\end{itemize}

\subsection{Diszjunkció (∨)}
Az A ∨ B jelentése, hogy legalább az egyik (A vagy B) igaz.
\begin{code}%
%
\>[6]\AgdaOperator{\AgdaField{\AgdaUnderscore{}∨\AgdaUnderscore{}}}%
\>[12]\AgdaSymbol{:}\AgdaSpace{}%
\AgdaSymbol{∀\{}\AgdaBound{Γ}\AgdaSymbol{\}}\AgdaSpace{}%
\AgdaSymbol{→}\AgdaSpace{}%
\AgdaField{For}\AgdaSpace{}%
\AgdaBound{Γ}\AgdaSpace{}%
\AgdaSymbol{→}\AgdaSpace{}%
\AgdaField{For}\AgdaSpace{}%
\AgdaBound{Γ}\AgdaSpace{}%
\AgdaSymbol{→}\AgdaSpace{}%
\AgdaField{For}\AgdaSpace{}%
\AgdaBound{Γ}\<%
\end{code}
\begin{code}%
%
\>[6]\AgdaField{∨[]}%
\>[12]\AgdaSymbol{:}\AgdaSpace{}%
\AgdaSymbol{∀\{}\AgdaBound{Γ}\AgdaSpace{}%
\AgdaBound{A}\AgdaSpace{}%
\AgdaBound{B}\AgdaSpace{}%
\AgdaBound{Δ}\AgdaSymbol{\}\{}\AgdaBound{γ}\AgdaSpace{}%
\AgdaSymbol{:}\AgdaSpace{}%
\AgdaField{Sub}\AgdaSpace{}%
\AgdaBound{Δ}\AgdaSpace{}%
\AgdaBound{Γ}\AgdaSymbol{\}}\AgdaSpace{}%
\AgdaSymbol{→}\AgdaSpace{}%
\AgdaSymbol{(}\AgdaBound{A}\AgdaSpace{}%
\AgdaOperator{\AgdaField{∨}}\AgdaSpace{}%
\AgdaBound{B}\AgdaSymbol{)}\AgdaSpace{}%
\AgdaOperator{\AgdaField{[}}\AgdaSpace{}%
\AgdaBound{γ}\AgdaSpace{}%
\AgdaOperator{\AgdaField{]ᶠ}}\AgdaSpace{}%
\AgdaOperator{\AgdaDatatype{≡}}\AgdaSpace{}%
\AgdaBound{A}\AgdaSpace{}%
\AgdaOperator{\AgdaField{[}}\AgdaSpace{}%
\AgdaBound{γ}\AgdaSpace{}%
\AgdaOperator{\AgdaField{]ᶠ}}\AgdaSpace{}%
\AgdaOperator{\AgdaField{∨}}\AgdaSpace{}%
\AgdaBound{B}\AgdaSpace{}%
\AgdaOperator{\AgdaField{[}}\AgdaSpace{}%
\AgdaBound{γ}\AgdaSpace{}%
\AgdaOperator{\AgdaField{]ᶠ}}\<%
\end{code}
\begin{itemize}
\item A diszjunkció bevezető szabálya (\AgdaField{∨in₁}): Ha A igaz, akkor A ∨ B is igaz.
\item A diszjunkció bevezető szabálya (\AgdaField{∨in₂}): Ha B igaz, akkor A ∨ B is igaz.
\begin{code}%
%
\>[6]\AgdaField{∨in₁}%
\>[12]\AgdaSymbol{:}\AgdaSpace{}%
\AgdaSymbol{∀\{}\AgdaBound{Γ}\AgdaSpace{}%
\AgdaBound{A}\AgdaSpace{}%
\AgdaBound{B}\AgdaSymbol{\}}\AgdaSpace{}%
\AgdaSymbol{→}\AgdaSpace{}%
\AgdaField{Pf}\AgdaSpace{}%
\AgdaBound{Γ}\AgdaSpace{}%
\AgdaBound{A}\AgdaSpace{}%
\AgdaSymbol{→}\AgdaSpace{}%
\AgdaField{Pf}\AgdaSpace{}%
\AgdaBound{Γ}\AgdaSpace{}%
\AgdaSymbol{(}\AgdaBound{A}\AgdaSpace{}%
\AgdaOperator{\AgdaField{∨}}\AgdaSpace{}%
\AgdaBound{B}\AgdaSymbol{)}\<%
\\
%
\>[6]\AgdaField{∨in₂}%
\>[12]\AgdaSymbol{:}\AgdaSpace{}%
\AgdaSymbol{∀\{}\AgdaBound{Γ}\AgdaSpace{}%
\AgdaBound{A}\AgdaSpace{}%
\AgdaBound{B}\AgdaSymbol{\}}\AgdaSpace{}%
\AgdaSymbol{→}\AgdaSpace{}%
\AgdaField{Pf}\AgdaSpace{}%
\AgdaBound{Γ}\AgdaSpace{}%
\AgdaBound{B}\AgdaSpace{}%
\AgdaSymbol{→}\AgdaSpace{}%
\AgdaField{Pf}\AgdaSpace{}%
\AgdaBound{Γ}\AgdaSpace{}%
\AgdaSymbol{(}\AgdaBound{A}\AgdaSpace{}%
\AgdaOperator{\AgdaField{∨}}\AgdaSpace{}%
\AgdaBound{B}\AgdaSymbol{)}\<%
\end{code}
\item A diszjunkció kivezető szabálya (\AgdaField{∨out}): Ha tudjuk, hogy A ∨ B igaz, és abból is következik, ha A igaz, de abból is, ha B igaz, akkor C igaz.
\begin{code}%
%
\>[6]\AgdaField{∨out}%
\>[12]\AgdaSymbol{:}%
\>[755I]\AgdaSymbol{∀\{}\AgdaBound{Γ}\AgdaSpace{}%
\AgdaBound{A}\AgdaSpace{}%
\AgdaBound{B}\AgdaSpace{}%
\AgdaBound{C}\AgdaSymbol{\}}\AgdaSpace{}%
\AgdaSymbol{→}\AgdaSpace{}%
\AgdaField{Pf}\AgdaSpace{}%
\AgdaSymbol{(}\AgdaBound{Γ}\AgdaSpace{}%
\AgdaOperator{\AgdaField{▹ₚ}}\AgdaSpace{}%
\AgdaBound{A}\AgdaSymbol{)}\AgdaSpace{}%
\AgdaSymbol{(}\AgdaBound{C}\AgdaSpace{}%
\AgdaOperator{\AgdaField{[}}\AgdaSpace{}%
\AgdaField{pₚ}\AgdaSpace{}%
\AgdaOperator{\AgdaField{]ᶠ}}\AgdaSymbol{)}\AgdaSpace{}%
\AgdaSymbol{→}\AgdaSpace{}%
\AgdaField{Pf}\AgdaSpace{}%
\AgdaSymbol{(}\AgdaBound{Γ}\AgdaSpace{}%
\AgdaOperator{\AgdaField{▹ₚ}}\AgdaSpace{}%
\AgdaBound{B}\AgdaSymbol{)}\AgdaSpace{}%
\AgdaSymbol{(}\AgdaBound{C}\AgdaSpace{}%
\AgdaOperator{\AgdaField{[}}\AgdaSpace{}%
\AgdaField{pₚ}\AgdaSpace{}%
\AgdaOperator{\AgdaField{]ᶠ}}\AgdaSymbol{)}\<%
\\
\>[.][@{}l@{}]\<[755I]%
\>[14]\AgdaSymbol{→}\AgdaSpace{}%
\AgdaField{Pf}\AgdaSpace{}%
\AgdaBound{Γ}\AgdaSpace{}%
\AgdaSymbol{(}\AgdaBound{A}\AgdaSpace{}%
\AgdaOperator{\AgdaField{∨}}\AgdaSpace{}%
\AgdaBound{B}\AgdaSymbol{)}\AgdaSpace{}%
\AgdaSymbol{→}\AgdaSpace{}%
\AgdaField{Pf}\AgdaSpace{}%
\AgdaBound{Γ}\AgdaSpace{}%
\AgdaBound{C}\<%
\end{code}
\end{itemize}

\subsection{Konstans igaz (⊤)}
A \AgdaField{⊤} egy mindig igaz állítást jelöl.
\begin{code}%
%
\>[6]\AgdaField{⊤}%
\>[12]\AgdaSymbol{:}\AgdaSpace{}%
\AgdaSymbol{∀\{}\AgdaBound{Γ}\AgdaSymbol{\}}\AgdaSpace{}%
\AgdaSymbol{→}\AgdaSpace{}%
\AgdaField{For}\AgdaSpace{}%
\AgdaBound{Γ}\<%
\end{code}
\begin{code}%
%
\>[6]\AgdaField{⊤[]}%
\>[12]\AgdaSymbol{:}\AgdaSpace{}%
\AgdaSymbol{∀\{}\AgdaBound{Γ}\AgdaSpace{}%
\AgdaBound{Δ}\AgdaSymbol{\}\{}\AgdaBound{γ}\AgdaSpace{}%
\AgdaSymbol{:}\AgdaSpace{}%
\AgdaField{Sub}\AgdaSpace{}%
\AgdaBound{Δ}\AgdaSpace{}%
\AgdaBound{Γ}\AgdaSymbol{\}}\AgdaSpace{}%
\AgdaSymbol{→}\AgdaSpace{}%
\AgdaField{⊤}\AgdaSpace{}%
\AgdaOperator{\AgdaField{[}}\AgdaSpace{}%
\AgdaBound{γ}\AgdaSpace{}%
\AgdaOperator{\AgdaField{]ᶠ}}\AgdaSpace{}%
\AgdaOperator{\AgdaDatatype{≡}}\AgdaSpace{}%
\AgdaField{⊤}\<%
\end{code}
\begin{itemize}
\item Az \textit{igaz} bevezető szabálya (\AgdaField{⊤in}): Minden kontextusban igaz.
\begin{code}%
%
\>[6]\AgdaField{⊤in}%
\>[12]\AgdaSymbol{:}\AgdaSpace{}%
\AgdaSymbol{∀\{}\AgdaBound{Γ}\AgdaSymbol{\}}\AgdaSpace{}%
\AgdaSymbol{→}\AgdaSpace{}%
\AgdaField{Pf}\AgdaSpace{}%
\AgdaBound{Γ}\AgdaSpace{}%
\AgdaField{⊤}\<%
\end{code}
\end{itemize}

\subsection{Konstans hamis (⊥)}
Az \AgdaField{⊥} egy ellentmondást vagy hamis állítást jelöl.
\begin{code}%
%
\>[6]\AgdaField{⊥}%
\>[12]\AgdaSymbol{:}\AgdaSpace{}%
\AgdaSymbol{∀\{}\AgdaBound{Γ}\AgdaSymbol{\}}\AgdaSpace{}%
\AgdaSymbol{→}\AgdaSpace{}%
\AgdaField{For}\AgdaSpace{}%
\AgdaBound{Γ}\<%
\end{code}
\begin{code}%
%
\>[6]\AgdaField{⊥[]}%
\>[12]\AgdaSymbol{:}\AgdaSpace{}%
\AgdaSymbol{∀\{}\AgdaBound{Γ}\AgdaSpace{}%
\AgdaBound{Δ}\AgdaSymbol{\}\{}\AgdaBound{γ}\AgdaSpace{}%
\AgdaSymbol{:}\AgdaSpace{}%
\AgdaField{Sub}\AgdaSpace{}%
\AgdaBound{Δ}\AgdaSpace{}%
\AgdaBound{Γ}\AgdaSymbol{\}}\AgdaSpace{}%
\AgdaSymbol{→}\AgdaSpace{}%
\AgdaField{⊥}\AgdaSpace{}%
\AgdaOperator{\AgdaField{[}}\AgdaSpace{}%
\AgdaBound{γ}\AgdaSpace{}%
\AgdaOperator{\AgdaField{]ᶠ}}\AgdaSpace{}%
\AgdaOperator{\AgdaDatatype{≡}}\AgdaSpace{}%
\AgdaField{⊥}\<%
\end{code}
\begin{itemize}
\item A \textit{hamis} kivezető szabálya (\AgdaField{⊥out}): Ha egy bizonyítás során elérünk egy ellentmondást (\AgdaField{⊥}), akkor bármilyen tetszőleges állítást igaznak vehetünk (\textit{ex falso quodlibet}).
\begin{code}%
%
\>[6]\AgdaField{⊥out}%
\>[12]\AgdaSymbol{:}\AgdaSpace{}%
\AgdaSymbol{∀\{}\AgdaBound{Γ}\AgdaSpace{}%
\AgdaBound{A}\AgdaSymbol{\}}\AgdaSpace{}%
\AgdaSymbol{→}\AgdaSpace{}%
\AgdaField{Pf}\AgdaSpace{}%
\AgdaBound{Γ}\AgdaSpace{}%
\AgdaField{⊥}\AgdaSpace{}%
\AgdaSymbol{→}\AgdaSpace{}%
\AgdaField{Pf}\AgdaSpace{}%
\AgdaBound{Γ}\AgdaSpace{}%
\AgdaBound{A}\<%
\end{code}
\end{itemize}

\subsection{Kvantorok}

A kvantorok formális logikai eszközök, amelyek segítségével általános és létezési állításokat tehetünk. Bindernek is nevezzük őket, utalva rá, hogy egy változót köt, azaz meghatározza, hogy egy változó miként értelmezhető egy adott hatókörben.

\subsubsection{Univerzális kvantor (∀)}

Az univerzális kvantor segítségével tudjuk kifejezni, hogy egy állítás minden lehetséges elemre igaz. Például a $∀x P(x)$ azt jelenti, hogy minden x-re igaz, hogy $P(x)$. A deklarációban látható is, hogy mivel egy binder-ről beszélünk, egy formulában lévő szabad változót kötünk meg, így egy szükebb kontextusban lévő formulát fogunk kapni.
\begin{code}%
%
\>[6]\AgdaField{Forall}\AgdaSpace{}%
\AgdaSymbol{:}\AgdaSpace{}%
\AgdaSymbol{∀\{}\AgdaBound{Γ}\AgdaSymbol{\}}\AgdaSpace{}%
\AgdaSymbol{→}\AgdaSpace{}%
\AgdaField{For}\AgdaSpace{}%
\AgdaSymbol{(}\AgdaBound{Γ}\AgdaSpace{}%
\AgdaOperator{\AgdaField{▹ₜ}}\AgdaSymbol{)}\AgdaSpace{}%
\AgdaSymbol{→}\AgdaSpace{}%
\AgdaField{For}\AgdaSpace{}%
\AgdaBound{Γ}\<%
\end{code}
\begin{code}%
%
\>[6]\AgdaField{Forall[]}\AgdaSpace{}%
\AgdaSymbol{:}%
\>[845I]\AgdaSymbol{∀\{}\AgdaBound{Γ}\AgdaSpace{}%
\AgdaBound{A}\AgdaSpace{}%
\AgdaBound{Δ}\AgdaSymbol{\}\{}\AgdaBound{γ}\AgdaSpace{}%
\AgdaSymbol{:}\AgdaSpace{}%
\AgdaField{Sub}\AgdaSpace{}%
\AgdaBound{Δ}\AgdaSpace{}%
\AgdaBound{Γ}\AgdaSymbol{\}}\<%
\\
\>[.][@{}l@{}]\<[845I]%
\>[17]\AgdaSymbol{→}\AgdaSpace{}%
\AgdaField{Forall}\AgdaSpace{}%
\AgdaBound{A}\AgdaSpace{}%
\AgdaOperator{\AgdaField{[}}\AgdaSpace{}%
\AgdaBound{γ}\AgdaSpace{}%
\AgdaOperator{\AgdaField{]ᶠ}}\AgdaSpace{}%
\AgdaOperator{\AgdaDatatype{≡}}\AgdaSpace{}%
\AgdaField{Forall}\AgdaSpace{}%
\AgdaSymbol{(}\AgdaBound{A}\AgdaSpace{}%
\AgdaOperator{\AgdaField{[}}\AgdaSpace{}%
\AgdaBound{γ}\AgdaSpace{}%
\AgdaOperator{\AgdaField{∘}}\AgdaSpace{}%
\AgdaField{pₜ}\AgdaSpace{}%
\AgdaOperator{\AgdaField{,ₜ}}\AgdaSpace{}%
\AgdaField{qₜ}\AgdaSpace{}%
\AgdaOperator{\AgdaField{]ᶠ}}\AgdaSymbol{)}\<%
\end{code}
\begin{itemize}
\item Az univerzális kvantor bevezető szabálya (\AgdaField{∀in}):
\begin{code}%
%
\>[6]\AgdaField{∀in}\AgdaSpace{}%
\AgdaSymbol{:}\AgdaSpace{}%
\AgdaSymbol{∀\{}\AgdaBound{Γ}\AgdaSpace{}%
\AgdaBound{A}\AgdaSymbol{\}}\AgdaSpace{}%
\AgdaSymbol{→}\AgdaSpace{}%
\AgdaField{Pf}\AgdaSpace{}%
\AgdaSymbol{(}\AgdaBound{Γ}\AgdaSpace{}%
\AgdaOperator{\AgdaField{▹ₜ}}\AgdaSpace{}%
\AgdaSymbol{)}\AgdaSpace{}%
\AgdaBound{A}\AgdaSpace{}%
\AgdaSymbol{→}\AgdaSpace{}%
\AgdaField{Pf}\AgdaSpace{}%
\AgdaBound{Γ}\AgdaSpace{}%
\AgdaSymbol{(}\AgdaField{Forall}\AgdaSpace{}%
\AgdaBound{A}\AgdaSymbol{)}\<%
\end{code}
\item Az univerzális kvantor kivezető szabálya (\AgdaField{∀out}):
\begin{code}%
%
\>[6]\AgdaField{∀out}\AgdaSpace{}%
\AgdaSymbol{:}\AgdaSpace{}%
\AgdaSymbol{∀\{}\AgdaBound{Γ}\AgdaSpace{}%
\AgdaBound{A}\AgdaSymbol{\}}\AgdaSpace{}%
\AgdaSymbol{→}\AgdaSpace{}%
\AgdaField{Pf}\AgdaSpace{}%
\AgdaBound{Γ}\AgdaSpace{}%
\AgdaSymbol{(}\AgdaField{Forall}\AgdaSpace{}%
\AgdaBound{A}\AgdaSymbol{)}\AgdaSpace{}%
\AgdaSymbol{→}\AgdaSpace{}%
\AgdaField{Pf}\AgdaSpace{}%
\AgdaSymbol{(}\AgdaBound{Γ}\AgdaSpace{}%
\AgdaOperator{\AgdaField{▹ₜ}}\AgdaSpace{}%
\AgdaSymbol{)}\AgdaSpace{}%
\AgdaBound{A}\<%
\end{code}
\end{itemize}

\subsubsection{Egzisztenciális kvantor (∃)}

Az egzisztenciális kvantor azt fejezi ki, hogy létezik legalább egy olyan elem, amelyre igaz egy állítás. Például a $∃x P(x)$ azt jelenti, hogy létezik legalább egy x, amelyre igaz, hogy $P(x)$.
\begin{code}%
%
\>[6]\AgdaField{∃}\AgdaSpace{}%
\AgdaSymbol{:}\AgdaSpace{}%
\AgdaSymbol{∀\{}\AgdaBound{Γ}\AgdaSymbol{\}}\AgdaSpace{}%
\AgdaSymbol{→}\AgdaSpace{}%
\AgdaField{For}\AgdaSpace{}%
\AgdaSymbol{(}\AgdaBound{Γ}\AgdaSpace{}%
\AgdaOperator{\AgdaField{▹ₜ}}\AgdaSpace{}%
\AgdaSymbol{)}\AgdaSpace{}%
\AgdaSymbol{→}\AgdaSpace{}%
\AgdaField{For}\AgdaSpace{}%
\AgdaBound{Γ}\<%
\end{code}
\begin{code}%
%
\>[6]\AgdaField{∃[]}\AgdaSpace{}%
\AgdaSymbol{:}\AgdaSpace{}%
\AgdaSymbol{∀\{}\AgdaBound{Γ}\AgdaSpace{}%
\AgdaBound{A}\AgdaSpace{}%
\AgdaBound{Δ}\AgdaSymbol{\}\{}\AgdaBound{γ}\AgdaSpace{}%
\AgdaSymbol{:}\AgdaSpace{}%
\AgdaField{Sub}\AgdaSpace{}%
\AgdaBound{Δ}\AgdaSpace{}%
\AgdaBound{Γ}\AgdaSymbol{\}}\AgdaSpace{}%
\AgdaSymbol{→}\AgdaSpace{}%
\AgdaField{∃}\AgdaSpace{}%
\AgdaBound{A}\AgdaSpace{}%
\AgdaOperator{\AgdaField{[}}\AgdaSpace{}%
\AgdaBound{γ}\AgdaSpace{}%
\AgdaOperator{\AgdaField{]ᶠ}}\AgdaSpace{}%
\AgdaOperator{\AgdaDatatype{≡}}\AgdaSpace{}%
\AgdaField{∃}\AgdaSpace{}%
\AgdaSymbol{(}\AgdaBound{A}\AgdaSpace{}%
\AgdaOperator{\AgdaField{[}}\AgdaSpace{}%
\AgdaBound{γ}\AgdaSpace{}%
\AgdaOperator{\AgdaField{∘}}\AgdaSpace{}%
\AgdaField{pₜ}\AgdaSpace{}%
\AgdaOperator{\AgdaField{,ₜ}}\AgdaSpace{}%
\AgdaField{qₜ}\AgdaSpace{}%
\AgdaOperator{\AgdaField{]ᶠ}}\AgdaSymbol{)}\<%
\end{code}
\begin{itemize}
\item Az egzisztenciális kvantor bevezető szabálya (\AgdaField{∃in}):
\begin{code}%
%
\>[6]\AgdaField{∃in}\AgdaSpace{}%
\AgdaSymbol{:}\AgdaSpace{}%
\AgdaSymbol{∀\{}\AgdaBound{Γ}\AgdaSpace{}%
\AgdaBound{A}\AgdaSymbol{\}}\AgdaSpace{}%
\AgdaSymbol{→}\AgdaSpace{}%
\AgdaSymbol{(}\AgdaBound{t}\AgdaSpace{}%
\AgdaSymbol{:}\AgdaSpace{}%
\AgdaField{Tm}\AgdaSpace{}%
\AgdaBound{Γ}\AgdaSymbol{)}\AgdaSpace{}%
\AgdaSymbol{→}\AgdaSpace{}%
\AgdaField{Pf}\AgdaSpace{}%
\AgdaBound{Γ}\AgdaSpace{}%
\AgdaSymbol{(}\AgdaBound{A}\AgdaSpace{}%
\AgdaOperator{\AgdaField{[}}\AgdaSpace{}%
\AgdaField{id}\AgdaSpace{}%
\AgdaOperator{\AgdaField{,ₜ}}\AgdaSpace{}%
\AgdaBound{t}\AgdaSpace{}%
\AgdaOperator{\AgdaField{]ᶠ}}\AgdaSymbol{)}\AgdaSpace{}%
\AgdaSymbol{→}\AgdaSpace{}%
\AgdaField{Pf}\AgdaSpace{}%
\AgdaBound{Γ}\AgdaSpace{}%
\AgdaSymbol{(}\AgdaField{∃}\AgdaSpace{}%
\AgdaBound{A}\AgdaSymbol{)}\<%
\end{code}
\item Az egzisztenciális kvantor kivezető szabálya (\AgdaField{∃out}):
\begin{code}%
%
\>[6]\AgdaField{∃out}\AgdaSpace{}%
\AgdaSymbol{:}\AgdaSpace{}%
\AgdaSymbol{∀\{}\AgdaBound{Γ}\AgdaSpace{}%
\AgdaBound{A}\AgdaSpace{}%
\AgdaBound{C}\AgdaSymbol{\}}\AgdaSpace{}%
\AgdaSymbol{→}\AgdaSpace{}%
\AgdaField{Pf}\AgdaSpace{}%
\AgdaSymbol{(}\AgdaBound{Γ}\AgdaSpace{}%
\AgdaOperator{\AgdaField{▹ₜ}}\AgdaSpace{}%
\AgdaOperator{\AgdaField{▹ₚ}}\AgdaSpace{}%
\AgdaBound{A}\AgdaSymbol{)}\AgdaSpace{}%
\AgdaSymbol{(}\AgdaBound{C}\AgdaSpace{}%
\AgdaOperator{\AgdaField{[}}\AgdaSpace{}%
\AgdaField{pₜ}\AgdaSpace{}%
\AgdaOperator{\AgdaField{∘}}\AgdaSpace{}%
\AgdaField{pₚ}\AgdaSpace{}%
\AgdaOperator{\AgdaField{]ᶠ}}\AgdaSymbol{)}\AgdaSpace{}%
\AgdaSymbol{→}\AgdaSpace{}%
\AgdaField{Pf}\AgdaSpace{}%
\AgdaBound{Γ}\AgdaSpace{}%
\AgdaSymbol{(}\AgdaField{∃}\AgdaSpace{}%
\AgdaBound{A}\AgdaSymbol{)}\AgdaSpace{}%
\AgdaSymbol{→}\AgdaSpace{}%
\AgdaField{Pf}\AgdaSpace{}%
\AgdaBound{Γ}\AgdaSpace{}%
\AgdaBound{C}\<%
\end{code}
\end{itemize}

\begin{code}[hide]%
%
\>[3]\AgdaKeyword{infixl}\AgdaSpace{}%
\AgdaNumber{5}\AgdaSpace{}%
\AgdaOperator{\AgdaField{\AgdaUnderscore{}▹ₚ\AgdaUnderscore{}}}\<%
\\
%
\>[3]\AgdaKeyword{infixl}\AgdaSpace{}%
\AgdaNumber{5}\AgdaSpace{}%
\AgdaOperator{\AgdaField{\AgdaUnderscore{},ₚ\AgdaUnderscore{}}}\<%
\\
%
\>[3]\AgdaKeyword{infixl}\AgdaSpace{}%
\AgdaNumber{5}\AgdaSpace{}%
\AgdaOperator{\AgdaField{\AgdaUnderscore{}▹ₜ}}\<%
\\
%
\>[3]\AgdaKeyword{infixl}\AgdaSpace{}%
\AgdaNumber{5}\AgdaSpace{}%
\AgdaOperator{\AgdaField{\AgdaUnderscore{},ₜ\AgdaUnderscore{}}}\<%
\\
%
\>[3]\AgdaKeyword{infixr}\AgdaSpace{}%
\AgdaNumber{6}\AgdaSpace{}%
\AgdaOperator{\AgdaField{\AgdaUnderscore{}∘\AgdaUnderscore{}}}\<%
\\
%
\>[3]\AgdaKeyword{infixr}\AgdaSpace{}%
\AgdaNumber{6}\AgdaSpace{}%
\AgdaOperator{\AgdaField{\AgdaUnderscore{}⊃\AgdaUnderscore{}}}\<%
\\
%
\>[3]\AgdaKeyword{infixr}\AgdaSpace{}%
\AgdaNumber{7}\AgdaSpace{}%
\AgdaOperator{\AgdaField{\AgdaUnderscore{}∨\AgdaUnderscore{}}}\<%
\\
%
\>[3]\AgdaKeyword{infixr}\AgdaSpace{}%
\AgdaNumber{8}\AgdaSpace{}%
\AgdaOperator{\AgdaField{\AgdaUnderscore{}∧\AgdaUnderscore{}}}\<%
\\
%
\>[3]\AgdaKeyword{infixl}\AgdaSpace{}%
\AgdaNumber{9}\AgdaSpace{}%
\AgdaOperator{\AgdaField{\AgdaUnderscore{}[\AgdaUnderscore{}]ᶠ}}\<%
\\
%
\>[3]\AgdaKeyword{infixl}\AgdaSpace{}%
\AgdaNumber{9}\AgdaSpace{}%
\AgdaOperator{\AgdaField{\AgdaUnderscore{}[\AgdaUnderscore{}]ᵖ}}\<%
\\
%
\>[3]\AgdaKeyword{infixl}\AgdaSpace{}%
\AgdaNumber{9}\AgdaSpace{}%
\AgdaOperator{\AgdaField{\AgdaUnderscore{}[\AgdaUnderscore{}]ᵗ}}\<%
\end{code}